%! Author = taj
%! Date = 7/22/25

% Preamble
\documentclass[11pt]{article}

% Packages
\usepackage{amsmath}
\usepackage{graphicx}
\usepackage{float}

\title{4. Linearni in kvadratni zakon upora}
\author{Taj Šobot}

% Document
\begin{document}
    \maketitle
    \flushleft

    V tem eksperimentu smo preučevali givanje predmetov skozi tekočine in gibanje tekočine mimo predmetov.
    Kako vpliva viskoznost, oblika in hitrost tekočine na silo ki deluje na telo.

    \section{Meritve}

    \subsection{1. Linearni zakon upora}

    Kroglico smo spuščali skozi viskozno olje.
    Kroglica se je kmalu začela gibati enakomerno.


    Takrat velja:
    \[
        F_g - F_u - F_{vzg} = 0
    \]
    \begin{figure}[H]
        \centering
        \includegraphics[width=0.9\textwidth]{img}
        \caption{Skica sistema kroglice}
        \label{fig:1}
    \end{figure}

    Te sile so:
    \begin{align*}
        F_g = m*g \\
        F_u = 6 \pi \eta r v \\
        F_{vzg} = 4/3 \pi r_k^3 \rho_0 g \\
    \end{align*}
    Po vstavljanju in obračanju dobimo to zvezo:
    \[
        \eta = \frac{2r_k^2 g(\rho_k - \rho_0)}{9v_k}
    \]
    Končno hitrost pa dobimo iz meritev z enačbo $L = (0.62 \pm 0.01) m$:
    \[
        v_k = \frac{L}{t}
    \]



    Meritve so prikazane v spodnji tabeli (večja in mala kroglica):

    \begin{table}[H]
        \centering
        \begin{tabular}{|c|c|c|}
            \hline
            n & t [s] (večja) & t [s] (manjša)\\
            \hline
            1 & 0.70 & 0.83 \\
            2 & 0.68 & 0.83 \\
            3 & 0.68 & 0.83 \\
            4 & 0.70 & 0.86 \\
            5 & 0.70 & 0.83 \\
            6 & 0.70 & 0.84 \\
            \hline
        \end{tabular}
        \caption{tabela 1}
        \label{tab:1}
    \end{table}
    \[
        t_{err} = \pm 0.1 s
    \]
    Velikosti kroglic:
    \begin{align*}
        R_v = (0,600 \pm 0,002)cm \\
        R_m = (0,35 \pm 0,002)cm
    \end{align*}
    Mase kroglic:
    \begin{align*}
        m_v = (7,1 \pm 0,1)g \\
        m_m = (1,4 \pm 0,1)g
    \end{align*}

    \subsection{2.a Linearni zakon upora}
    Pri tem delu vaje smo merili hitrost zraka z pitotovo cevjo.
    Zrak smo z vetrovnikom pospeševali.
    Z u-cevko smo merili zastojni tlak.
    Pitotovo cev smo tudi premikali pravokotno na premikajoči zrak, tako smo dobili profil hitrosti.

    Meritve:

    \begin{table}[H]
        \centering
        \begin{tabular}{|c|c|c|}
            \hline
            n & $\Delta x [cm]$ & $\Delta p [Pa]$\\
            \hline
            1	& 0.5 & 12  \\
            2	& 0.9 & 21  \\
            3	& 1.3 & 31  \\
            4	& 1.9 & 45  \\
            5	& 2.4 & 57  \\
            6	& 2.4 & 57  \\
            7	& 2.6 & 62  \\
            8	& 2.6 & 62  \\
            9	& 2.5 & 59  \\
            10	& 1.2 & 28  \\
            11	& 0.6 & 14  \\
            12	& 0.2 & 5   \\
            \hline
        \end{tabular}
        \caption{tabela tlakov}
        \label{tab:2}
    \end{table}
    Napake:
    \begin{align*}
        x_{err} = \pm 0.1 cm \\
        p_{err} = \pm 2 Pa
    \end{align*}

    U cevko smo meli nastavljeno pod kotom $\alpha = (14 \pm 1) ^\circ$
    Tlak (v zgornji tabeli) smo izračunali po tej enačbi:
    \[
        \Delta p = \rho g x \sin(\alpha)
    \]

    \subsection{2.b Merjenje koeficienta upora}
    Na isti način kot pri prejšnji vaji smo postavili pred vetrovnik različne objekte in merili silo upora na njih.
    Kasneje pa še izračunati koeficient upora za vsak predmet.
    To smo dosegli z vzvodom in silomerom.
    Skica eksperimenta:

    \begin{figure}[H]
        \centering
        \includegraphics[width=0.9\textwidth]{img2}
        \caption{Skica eksperimenta}
        \label{fig:2}
    \end{figure}
    Tako je zrak ustvarjal navor, ki je:
    \[
        F_u r_2 = F_d r_1
    \]
    Dolžine so bile izmerjene:
    \begin{align*}
        r_2 = (19,9 \pm 0,1) cm \\
        r_1 = (21,2 \pm 2) cm
    \end{align*}
    Naše meritve so meritve sile $F_d$ v nasledni tabeli (3 različni predmeti):

    \begin{table}[H]
        \centering
        \begin{tabular}{|c|c|c|c|}
            \hline
            n & $F [N]$ (krogla) & $F [N]$ (ploscek) & F [N] (avto)\\
            \hline
            1	& 0.02	& 0.071	 & 0.011 \\
            2	& 0.018	& 0.069	 & 0.012 \\
            3	& 0.02	& 0.069	 & 0.013 \\
            4	& 0.018	& 0.065	 & 0.012 \\
            5	& 0.021	& 0.067	 & 0.014 \\

            \hline
        \end{tabular}
        \caption{tabela sil}
        \label{tab:3}
    \end{table}
    Vsaki od predmetov smo morali tudi izmerit prečni presek:
    \begin{align*}
        r_{krog} = (2,5 \pm 0,1) cm \\
        r_{plo} = (2,5 \pm 0,1) cm \\
        a_{avt} = (4,0 \pm 0,1) cm \\
        b_{avt} = (3,0 \pm 0,1) cm \\
    \end{align*}

    \section{Rezultati in grafi}

    \subsection{1. Linearni zakon upora}
    %todo
    Končne hitrosti kroglic:
    \begin{align*}
        v_{k1} = (0.89 \pm 0.05) m/s \\
        v_{k2} = (0.74 \ pm 0.05) m/s
    \end{align*}
    Gostota kroglic in olja je:
    \begin{align*}
        \rho_{k1} = (7800 \pm 10) kg/m^3 \\
        \rho_{k2} = (7800 \ pm 10) kg/m^3 \\
        \rho_{olja} = (830 \ pm 10) kg/m^3 \\
    \end{align*}
    Vstavimo v enačbo:
    \[
        \eta = \frac{2r_k^2 g(\rho_k - \rho_0)}{9v_k}
    \]
    Dobimo:
    \[
        \eta_1 = (0.61 \pm 0.02) kg/ms \\
        \eta_2 = (0.58 \pm 0.02) kg/ms
    \]
    te vrednosti se strinjajo z realnimi.

    \subsection{2.a Merjenje hitrosti zraka}

    Hitrost zraka lahko izračunamo z enačbo:
    \[
        v = \sqrt{\frac{2 \rho g l \sin(\alpha)}{\rho_z}}
    \]

    Izracuni v tabeli:
    \begin{table}[H]
        \centering
        \begin{tabular}{|c|c|c|}
            \hline
            n & $v [m/s]$ \\
            \hline
            1	& 4.3  \\
            2	& 5.8  \\
            3	& 6.9  \\
            4	& 8.4  \\
            5	& 9.4  \\
            6	& 9.4  \\
            7	& 9.8  \\
            8	& 9.8  \\
            9	& 9.6  \\
            10	& 6.6  \\
            11	& 4.7  \\
            12	& 2.7  \\
            \hline
        \end{tabular}
        \caption{tabela hitrosti}
        \label{tab:4}
    \end{table}
    Kot vidimo je v sredini hitrost zraka okoli $v_z = (9.7 \pm 0.2) m/s$.

    \subsection{2.b Merjenje koeficienta upora}
    Koefecient upora za vsak predmet se izračuna z ($\rho_z = (1,29 \pm 0,02) kg/m^3$):
    \[
        C = \frac{2 F_u}{S \rho_z v^2}
    \]
    Izračunane povprečne meritve:
    \begin{align*}
        F_{krog} = (0.019 \pm 0.002) N \\
        F_{plo} = (0.034 \pm 0.003) N \\
        F_{avt} = (0.012 \pm 0.002) N
    \end{align*}

    Izračunane povprečne sile:
    \begin{align*}
        F_{ukrog} = (0.020 \pm 0.002) N \\
        F_{uplo} = (0.036 \pm 0.003) N \\
        F_{uavt} = (0.013 \pm 0.002) N
    \end{align*}
    Izračunane ploščine:
    \begin{align*}
        S_{krog} = (0.02 \pm 0.0003) m^2 \\
        S_{plo} = (0.003 \pm 0.0001) m^2  \\
        S_{avt} = (0.0012 \pm 0.0002) m^2
    \end{align*}
    Izračunani koeficienti:
    \begin{align*}
        C_{krog} = (0.23 \pm 0.05)  \\
        C_{plo} = (0.43 \pm 0.05)   \\
        C_{avt} = (0.26 \pm 0.05)
    \end{align*}
    Komentar: te vrednosti so malo manjše od realnih.


\end{document}